\documentclass[11pt]{article}
\usepackage{mathblatt}
% Präambel des Setzers (werkzeuge/setzer-praeambel.tex, wortgleich); dazu \gteil (Muster T2U)
\makeatletter
\setlength{\mbpfbe}{0mm}% keine Punkte (pruefung.md § 5: nur beim Prüfstein)
\newlength{\gEin}\setlength{\gEin}{0pt}
\let\gSchrift\relax
\renewcommand{\pfkreuzzeile}[1]{\par\noindent\hangindent3.2em\hangafter1\hspace*{1.6em}$\square$\ #1\par}
\newcommand{\gkreuz}[1]{$\square$\ #1\hspace{2em}}
\newcommand{\gkreuzzeile}[2]{\par\noindent\hangindent2.2em\hangafter1\mbox{#1\,$\square$}\ #2\par}% Buchstabe, Option: eine je Zeile (neben Skizze)
% links, was man ansieht, rechts, was man tut; Spaltengrenze überall an derselben Stelle (gemeinsam.md „Satz“)
\newsavebox{\gbox}\newlength{\gLinks}
\newcommand{\gzwei}[2]{\par\smallskip\noindent\sbox{\gbox}{#1}%
  \setlength{\gLinks}{\dimexpr0.5\textwidth-\gEin-\mbpfnr\relax}%
  \ifdim\wd\gbox>\dimexpr\gLinks-4mm\relax
    \usebox{\gbox}\par\smallskip\noindent\begin{minipage}{\linewidth}#2\end{minipage}%
  \else
    \begin{minipage}[t]{\gLinks}\vspace{0pt}\usebox{\gbox}\end{minipage}%
    \begin{minipage}[t]{\dimexpr\linewidth-\gLinks\relax}\vspace{0pt}\raggedright #2\end{minipage}%
  \fi\par}
\RenewDocumentEnvironment{pfaufg}{m m m}{%
  \begin{lrbox}{\mbaufgabebox}\begin{minipage}[t]{\linewidth}%
  \gSchrift\noindent\hspace*{\gEin}\mbpfrandmarke{#1}%
  \begin{minipage}[t]{\mbpfnr}\strut\textbf{#2}\end{minipage}%
  \begin{minipage}[t]{\dimexpr\linewidth-\gEin-\mbpfnr\relax}\raggedright\strut\ignorespaces}%
 {\end{minipage}%
  \end{minipage}\end{lrbox}%
  \setlength{\mbaufgabehoehe}{\dimexpr\ht\mbaufgabebox+\dp\mbaufgabebox\relax}%
  \par\addvspace{5pt}\Needspace*{\mbaufgabehoehe}%
  \noindent\usebox{\mbaufgabebox}\par\addvspace{5pt}}
% Teilaufgabe mit Skizze links in gzwei: Buchstabe in eigener Spalte, darunter Text und Skizze
\newcommand{\gteil}[2]{\makebox[\mbpfteil][l]{\textbf{#1}}\parbox[t]{\dimexpr0.5\textwidth-\gEin-\mbpfnr-5mm-\mbpfteil\relax}{\raggedright #2}}
% Stufenkopf ohne die festen 14 Zeilen Platz der Vorlage (Nachzug 12.10., Wunsch GYW): der Setzer
% fordert vor dem Kopf 0,16 Seitenhöhe (Kopf + Schülerfrage), mehr nicht
\RenewDocumentCommand{\pfstufekopf}{o m m}{\par\addvspace{10pt}%
  \noindent\IfValueT{#1}{\hypertarget{#1}{}\label{#1}}{\large\bfseries #2}%
  \ifblank{#3}{}{\hspace{0.8em}{\small\color{mbgrau}#3}}\par\nobreak\vspace{2pt}}
\makeatother
\newcommand{\karo}[1]{\par\smallskip\noindent\begin{tikzpicture}
  \clip (0,0) rectangle ({\linewidth-1mm},{#1*10mm});
  \draw[black!16,line width=0.3pt,step=5mm] (0,0) grid ({\linewidth},{#1*10mm});
  \end{tikzpicture}\par}
\newcommand{\kk}[1]{$\square$\,#1\hspace{0.9em}}
\newcommand{\linie}[1]{\,{\color{black!45}\rule[-1pt]{#1}{0.4pt}}\,}
\newcommand{\pkt}{\textperiodcentered{}}
% v0.7 (12.10.): Schülerfrage unter dem Stufenkopf; Ankreuzoption als Bild (Buchstabe und Kästchen darunter)
\newcommand{\gfrage}[1]{\par\nobreak\noindent\hspace*{\mbpfnr}{\small\itshape #1}\par\nobreak\addvspace{2pt}}
\newcommand{\gbildoption}[3]{\begin{minipage}[t][#3][b]{0.32\linewidth}\centering #2\par\smallskip\mbox{#1\,$\square$}\end{minipage}}
\begin{document}

\pfheftstil
\pflaufkopf[P10 \pkt{} Prüfung]{Fläche und Umfang \pkt{} Lösungen}{H37}\pfseitevon
\pfheftkopf{Fläche und Umfang}{Lösungen}
\pfstufenzeile{Formel zur Figur \pkt{} Fläche oder Umfang berechnen \pkt{} Seite aus Fläche oder Umfang \pkt{} Figur erst zerlegen \pkt{} Anteil in Prozent}

\Needspace*{0.16\textheight}\pfstufekopf{Formel zur Figur}{letzte 5 Jahre: 2025 \pkt{} 2024 \pkt{} 2022}
\begin{pfloesungb}{}
\lzb{1.}{}{a) A = a · b; b) A = g · h : 2; c) u = 2 · $\pi$ · r; d) A = (a + c) : 2 · h}{Fläche: zwei Längen mal; Umfang: einmal herum}
\end{pfloesungb}
\begin{pfloesungb}{}
\lzb{2.}{}{$6x$}{$u$; zweimal $2x$ und zweimal $x$}
\end{pfloesungb}
\begin{pfloesungb}{}
\lzb{3.}{}{C}{A = a · (b + c) (dritte Option)}
\end{pfloesungb}

\Needspace*{0.16\textheight}\pfstufekopf{Fläche oder Umfang berechnen}{letzte 5 Jahre: 2026 \pkt{} 2025 \pkt{} 2024 \pkt{} 2023}
\begin{pfloesungb}{}
\lzb{4.}{}{a) und c)}{Fläche bedeckt (Teppich, Farbe), Umfang läuft herum (Zaun, Bordüre)}
\end{pfloesungb}
\begin{pfloesungb}{}
\lzb{5.}{}{$22$ cm}{$A = 8 \cdot 3 = 24$ cm², $u = 2 \cdot (8 + 3)$}
\end{pfloesungb}
\begin{pfloesungb}{}
\lzb{6.}{}{14 cm}{14 cm}
\end{pfloesungb}
\begin{pfloesungb}{}
\lzb{7.}{}{A $\approx$ 1963,5 cm²}{A = $\pi$ · 25² $\approx$ 1963,5 cm²}
\end{pfloesungb}

\Needspace*{0.16\textheight}\pfstufekopf{Seite aus Fläche oder Umfang}{letzte 5 Jahre: 2023}
\begin{pfloesungb}{}
\lzb{8.}{}{a) b = A : a; b) a = √A; c) b = u : 2 $-$ a}{erst die Formel, dann umstellen}
\end{pfloesungb}
\begin{pfloesungb}{}
\lzb{9.}{}{$7$ cm}{$b = 56 : 8$}
\end{pfloesungb}
\begin{pfloesungb}{}
\lzb{10.}{}{110 mm}{110 mm}
\end{pfloesungb}

\Needspace*{0.16\textheight}\pfstufekopf{Figur erst zerlegen}{letzte 5 Jahre: 2023 \pkt{} 2022}
\begin{pfloesungb}{}
\lzb{11.}{}{a) Rechteck + Kreis; b) Quadrat + Dreieck; c) Rechteck $-$ Kreis}{zusammensetzen heißt plus, ein Loch heißt minus}
\end{pfloesungb}
\begin{pfloesungb}{}
\lzb{12.}{}{$36$ m²}{$A = 8 \cdot 3 + 3 \cdot 4 = 24 + 12$}
\end{pfloesungb}
\begin{pfloesungb}{}
\lzb{13.}{}{u $\approx$ 57,4 m}{Schenkel = 7,50 : sin 62° $\approx$ 8,49 m; u = 24,20 + 16,20 + 2 · Schenkel $\approx$ 57,4 m}
\end{pfloesungb}

\Needspace*{0.16\textheight}\pfstufekopf{Anteil in Prozent}{letzte 5 Jahre: 2025 \pkt{} 2023}
\begin{pfloesungb}{}
\lzb{14.}{}{a) 360°; b) die ganze Platte; c) das ganze Grundstück}{Anteil = Teil : Ganzes, das Ganze ist 100 \%}
\end{pfloesungb}
\begin{pfloesungb}{}
\lzb{15.}{}{Anteil grau $\approx$ 34,7 \%}{Anteil grau = 125 : 360 $\approx$ 34,7 \%}
\end{pfloesungb}
\end{document}